% from ega4-2-fr.tex
\oldpage[IV-2]{121}

premiers associés de codimension 1, l'hypothèse
$\dim(\mathcal O_{X,x})=1$ entraîne $x\notin\operatorname{Ass}(\mathcal O_X)$,
donc $\operatorname{prof}(\mathcal O_{X,x})=1$, et par suite aussi
$\operatorname{prof}(\mathcal O_{X',x'})=1$. Supposons maintenant
$\dim(\mathcal O_{X',x'})\geq2$, donc $\dim(\mathcal O_{X,x})\geq2$,
c'est-à-dire $x\in Z^{(2)}$; on en déduit que
$\operatorname{prof}(\mathcal O_{X',x'})\geq2$ par (5.10.12, (ii)). Ceci
établit les assertions de (iii).

Pour prouver (iv) il suffit de remplacer $\mathcal O_X$ par $\mathcal F$ dans
le raisonnement précédent, qui établit que $\mathcal F'$ vérifie $(S_2)$ et
que si $\dim(\mathcal F'_{x'})\leq1$, $\mathcal F_x$ et
$\mathcal F'_{x'}$ sont di-isomorphes; en particulier si
$\dim(\mathcal F'_{x'})=0$, on a $\dim(\mathcal F_x)=0$, donc
$\dim(\mathcal O_{X,x})=0$ puisque $\mathcal F$ a pour support $X$, et enfin
$\dim(\mathcal O_{X',x'})=0$; toute composante irréductible de
$\operatorname{Supp}(\mathcal F')$ est donc une composante irréductible de
$X'$, puisque $\mathcal F'$ est cohérent, donc
$\operatorname{Supp}(\mathcal F')$ fermé.

\begin{proposition}[5.10.17]
\label{IV.2.s5.10.17-fr}
Soient $A$ un anneau intègre noethérien, et désignons par $A^{(1)}$
l'intersection des anneaux locaux $A_{\mathfrak p}$, où $\mathfrak p$
parcourt l'ensemble des idéaux premiers de $A$ de hauteur 1. Supposons que
$A^{(1)}$ soit une $A$-algèbre finie. Alors:
\begin{enumerate}
\item[(i)] L'anneau $A^{(1)}$ vérifie la condition $(S_2)$.
\item[(ii)] L'ensemble $U$ des $\mathfrak p\in\operatorname{Spec}(A)$ tels
que l'homomorphisme canonique
$A_{\mathfrak p}\to(A^{(1)})_{\mathfrak p}$ soit bijectif est égal à
l'ensemble des $\mathfrak p$ tels que $A_{\mathfrak p}$ vérifie $(S_2)$; $U$
est ouvert dans $X=\operatorname{Spec}(A)$ et l'on a
$\operatorname{codim}(X-U,X)\geq2$.
\item[(iii)] Pour toute partie multiplicative $S$ de $A$, $(S^{-1}A)^{(1)}$
est une $(S^{-1}A)$-algèbre finie.
\item[(iv)] Soit $B$ une $A$-algèbre finie, intègre et contenant $A$. Alors
$B^{(1)}$ est une $B$-algèbre finie. En outre, pour tout idéal premier
$\mathfrak q$ de $B$, de hauteur 1, l'idéal premier $\mathfrak q\cap A$ de
$A$ est de hauteur 1.
\end{enumerate}
\end{proposition}

Si l'on tient compte de la formule (5.9.3.1), on voit que
$X'=\operatorname{Spec}(A^{(1)})$ est la $Z^{(2)}$-clôture de
$X=\operatorname{Spec}(A)$; comme $A$ n'a pas d'idéaux premiers associés
immergés, les propriétés (i) et (ii) sont des cas particuliers de (5.10.16,
(i), (ii) et (iii)). Pour prouver (iii), il suffit de remarquer que l'on a
$(S^{-1}A)^{(1)}=S^{-1}A^{(1)}$, ce qui est un cas particulier de (5.9.4): en
effet $S^{-1}A$ est un $A$-module plat, les idéaux premiers de $S^{-1}A$ sont
les idéaux $S^{-1}\mathfrak p$, où $\mathfrak p\in\operatorname{Spec}(A)$ ne
rencontre pas $S$, et l'on a
$\operatorname{ht}(S^{-1}\mathfrak p)=\operatorname{ht}(\mathfrak p)$. Comme
$A^{(1)}$ est une $A$-algèbre finie, $S^{-1}A^{(1)}$ est une
$S^{-1}A$-algèbre finie, d'où (iii).

Pour démontrer (iv), posons $Y=\operatorname{Spec}(B)$, et soit $f:Y\to X$ le
morphisme structural; comme il est fini, il résulte de (5.4.1) que pour tout
$y\in Y$, on a
$\dim(\mathcal O_{Y,y})\leq\dim(\mathcal O_{X,f(y)})$; donc, si
$T=f^{-1}(Z^{(2)}(X))$, on a $T\supset Z^{(2)}(Y)$. Montrons que
\[
\mathcal G=\mathcal H^0_{X/Z^{(2)}(X)}(f_*(\mathcal O_Y))
\]
est cohérent; en effet $f_*(\mathcal O_Y)=\widetilde B$, $B$ étant considéré
comme $A$-module; mais comme $B$ est une $A$-algèbre intègre finie, son corps
des fractions est fini sur le corps des fractions de $A$, donc $B$ est contenu
dans un $A$-module libre de type fini, et par suite (5.9.2, (i)) $\mathcal G$
est un sous-$\mathcal O_X$-Module quasi-cohérent de
$(\mathcal H^0_{X/Z^{(2)}(X)}(\mathcal O_X))^n$ pour un $n$ convenable; ce
dernier étant cohérent par hypothèse, il en est de même de $\mathcal G$. Or,
il résulte de la définition (5.9.1.2) que $\mathcal G$ est isomorphe à
$f_*(\mathcal H^0_{Y/T}(\mathcal O_Y))$; cela prouve a fortiori que
$\mathcal H^0_{Y/T}(\mathcal O_Y)$ est un $\mathcal O_Y$-Module cohérent. Il
résulte alors de (5.10.10, (ii)), appliqué à $\mathcal O_Y$ et au point
générique de $Y$, que l'on a $\operatorname{codim}(T,Y)\geq2$, c'est-à-dire
$T\subset Z^{(2)}(Y)$, et finalement $T=Z^{(2)}(Y)$. Cela prouve les deux
assertions de (iv).

